\section{Intrinsic lifted-channel graphs and analytic table rigidity}
\label{sec:v23-intrinsic-multichannel-rigidity}

The signed one-channel inverse is formulated in a channel-adapted oriented
frame.  In a periodic billiard a measured channel also carries a deck
translation.  A global inverse theorem must therefore reconstruct both the
relative placement of distinct channel frames and their placement relative to
the periodic lattice.

Let \(\Lambda_{\rm abs}\simeq\mathbb Z^2\) be a marked abstract lattice with
its fixed Euclidean Gram form and the orientation induced by its marked basis.
An oriented Euclidean realization is an orientation-preserving linear isometry
\[
 \iota:\Lambda_{\rm abs}\otimes\mathbb R\longrightarrow\mathbb R^2.
\]
The quotient obstacle orbits are labelled \(1,\ldots,N\).  An oriented
measured channel is indexed exactly as in Section~\ref{sec:v18-formal-setup}:
\[
                         e=(a_e,b_e,\ell_e),
 \qquad \ell_e\in\Lambda_{\rm abs},
\]
and joins a lift of obstacle \(a_e\) to the lift of obstacle \(b_e\) translated
by \(\ell_e\).  Its two incidences are denoted \((e,-)\) and \((e,+)\), with
\begin{equation}\label{eq:v23-incidence-labels}
 r(e,-)=a_e,\quad \nu(e,-)=0,
 \qquad
 r(e,+)=b_e,\quad \nu(e,+)=\ell_e.
\end{equation}
Reversing the channel replaces \((a,b,\ell)\) by
\((b,a,-\ell)\).

\begin{definition}[Intrinsic signed lifted-channel datum]
\label{def:v23-intrinsic-channel-datum}
For an oriented channel \(e=(a,b,\ell)\), its intrinsic datum consists of
\begin{enumerate}
\item the lifted channel label \((a,b,\ell)\) in the marked abstract lattice;
\item the physical onset \(g_e\);
\item the two same-type signed endpoint-law germs
\[
 \{\Pi^{e,\mathrm{end},\infty}_{00,d}:0<d<d_e\},\qquad
 \{\Pi^{e,\mathrm{end},\infty}_{11,d}:0<d<d_e\};
\]
\item the sign convention for the transverse endpoint coordinate.
\end{enumerate}
The law is recorded in the channel-adapted frame in which the first contact is
\((0,0)\), the second is \((g_e,0)\), and the oriented transverse axes are
those used by the signed endpoint coordinates.  No Euclidean placement of
this frame relative to the lattice or to any other measured channel is
supplied.
\end{definition}

Theorem~\ref{thm:v22-signed-rigidity} and
Lemma~\ref{lem:v23-analytic-germ-globalization} recover two complete oriented
analytic boundary images in the edge frame,
\[
                         C_{e,-},\qquad C_{e,+},
\]
representing the lifted boundaries at deck offsets \(0\) and \(\ell_e\).
The remaining inverse problem is a periodic Euclidean gluing problem.

\subsection{Analytic signatures and local matching}

Let \(C\) be a connected oriented real-analytic strictly convex closed curve,
parametrized by oriented arclength \(s\in\mathbb R/L\mathbb Z\), with signed
curvature \(\kappa_C(s)>0\).  At \(p=C(s_0)\) define
\[
 \mathfrak j_C(p)=
 \bigl(\kappa_C(s_0),\kappa_C'(s_0),\kappa_C''(s_0),\ldots\bigr).
\]
Let \(\mathrm{Sym}^+(C)\subset\mathrm{SE}(2)\) be the orientation-preserving
Euclidean symmetry group of the boundary image.

\begin{lemma}[Coincident analytic signatures are exactly symmetry orbits]
\label{lem:v23-signature-symmetry}
For framed points \(p,q\in C\),
\[
                         \mathfrak j_C(p)=\mathfrak j_C(q)
\]
if and only if an element of \(\mathrm{Sym}^+(C)\) carries the oriented framed
point \(p\) to the oriented framed point \(q\).
\end{lemma}

\begin{proof}
The reverse implication is immediate.  For the forward implication choose
arclength so that \(p=C(0)\) and \(q=C(a)\).  Equality of the complete
curvature jets gives equality of the real-analytic germs
\(\kappa_C(s)\) and \(\kappa_C(s+a)\), hence equality for every \(s\) by the
identity theorem.  Let \(A\in\mathrm{SE}(2)\) send the oriented point--tangent
frame at \(0\) to that at \(a\).  Uniqueness for the planar Frenet system gives
\(AC(s)=C(s+a)\) for all \(s\), so \(A\in\mathrm{Sym}^+(C)\).
\end{proof}

Thus a circle has a genuine continuous phase ambiguity and a cyclically
symmetric analytic obstacle may have a finite one.  These are actual inverse
ambiguities, not registration conventions.

\subsection{Admissible periodic gluings}

Let \(\tau_v(x)=x+v\).  An \emph{algebraic periodic gluing} of the intrinsic
data is a pair
\begin{equation}\label{eq:v23-periodic-gluing}
 \left(\iota,(A_e)_{e\in E}\right),
 \qquad
 A_e\in\mathrm{SE}(2),
\end{equation}
where \(\iota\) is an oriented Euclidean realization of
\(\Lambda_{\rm abs}\), such that whenever two incidences
\((e,\sigma)\), \((f,\sigma')\) carry the same quotient obstacle label,
\begin{equation}\label{eq:v23-periodic-gluing-constraint}
 \tau_{-\iota(\nu(e,\sigma))}A_eC_{e,\sigma}
 =
 \tau_{-\iota(\nu(f,\sigma'))}A_fC_{f,\sigma'}
\end{equation}
as oriented labelled boundary images.  Thus every recovered lift is translated
back by its deck vector before copies of the same quotient obstacle are
identified.

The common values in \eqref{eq:v23-periodic-gluing-constraint} define quotient
obstacle images \(C_1,\ldots,C_N\).  Extend them by all deck translations
\(\iota(\Lambda_{\rm abs})\).  We call the gluing \emph{admissible} if this
periodic family is a dispersing billiard table: distinct translated obstacle
closures are disjoint and every measured edge has the prescribed facing
contacts, onset and selected closest-pair channel geometry.  These are
conditions on the reconstructed curves and their candidate placements.

The group \(\mathrm{SE}(2)\) acts by
\[
 B=(R,t):\qquad
 \iota\mapsto R\iota,
 \qquad A_e\mapsto B\circ A_e.
\]
Both \eqref{eq:v23-periodic-gluing-constraint} and admissibility are invariant.
Let \(\mathfrak G_{\rm per}(\mathscr D)\) be the quotient of the admissible
periodic gluings by this action.

\begin{theorem}[Intrinsic periodic reconstruction and holonomy classification]
\label{thm:v23-intrinsic-table-rigidity}
Let two labelled periodic dispersing billiard tables have the same marked
abstract Euclidean lattice, the same obstacle-orbit labels, and the same
connected measured lifted-channel graph.  Assume every obstacle boundary is
connected, real analytic and strictly convex, and that the two tables have
identical intrinsic signed lifted-channel data for every measured edge.  Then:
\begin{enumerate}
\item every incident lifted obstacle boundary is recovered as a complete
analytic image in its own channel frame;
\item the observations determine the admissible periodic gluing space
\(\mathfrak G_{\rm per}(\mathscr D)\) without a common laboratory registration;
\item labelled periodic dispersing tables realizing the observations, modulo
one simultaneous orientation-preserving Euclidean motion of the lattice and
all obstacle images, are in bijection with
\(\mathfrak G_{\rm per}(\mathscr D)\).
\end{enumerate}
In particular, if \(\mathfrak G_{\rm per}(\mathscr D)\) is a singleton, the
complete labelled periodic table is intrinsically determined up to one global
\(\mathrm{SE}(2)\) motion.

For the common-orientation orbit of the entire marked datum, as defined in
Definition~\ref{def:v28-common-orientation-datum}, the corresponding
classification is modulo \(\mathrm E(2)\), by
Theorem~\ref{thm:v28-unoriented-classification}.  This assertion concerns one
common orientation choice, not pointwise deletion of endpoint signs.
\end{theorem}

\begin{proof}
The first assertion follows from Theorem~\ref{thm:v22-signed-rigidity} and
Lemma~\ref{lem:v23-analytic-germ-globalization}.  Hence every curve entering
\eqref{eq:v23-periodic-gluing-constraint} is determined by the observations.

A table realization supplies its lattice realization \(\iota\) and the actual
placement \(A_e\) of every channel frame.  The incidence \((e,\sigma)\) lies
on the lift
\(D_{r(e,\sigma)}+\iota(\nu(e,\sigma))\), so translating it back by the deck
vector gives the same quotient obstacle image at every incidence with that
label.  The resulting gluing is admissible, and simultaneous global Euclidean
motion changes only its quotient representative.

Conversely an admissible gluing defines
\[
 C_r=\tau_{-\iota(\nu(e,\sigma))}A_eC_{e,\sigma}
\]
for any incidence with \(r(e,\sigma)=r\).  The gluing constraint makes this
independent of the choice.  The periodic lifts are
\(C_r+\iota(\lambda)\), \(\lambda\in\Lambda_{\rm abs}\).  For a measured edge
\(e=(a,b,\ell)\), its placed endpoint curves are exactly
\(C_a\) and \(C_b+\iota(\ell)\), with the observed onset and oriented signed
contact convention.  Admissibility makes the full periodic family a dispersing
billiard table with the declared measured channels.  These constructions are
inverse modulo the common \(\mathrm{SE}(2)\) action.
\end{proof}

\subsection{Cycle holonomy and an intrinsic uniqueness criterion}

If two incidences of the same obstacle orbit have contact signatures occurring
at unique framed points, Lemma~\ref{lem:v23-signature-symmetry} gives a unique
orientation-preserving Euclidean congruence between their recovered curve
copies.  We call such an incidence transition \emph{signature-rigid}.

Consider a closed oriented channel cycle
\[
 c=e_1^{\epsilon_1}\cdots e_m^{\epsilon_m},
 \qquad \epsilon_i\in\{\pm1\},
\]
for which every required incidence transition is signature-rigid.  Its marked
deck displacement is
\begin{equation}\label{eq:v23-cycle-deck}
 \eta_c=\sum_{i=1}^m\epsilon_i\ell_{e_i}
 \in\Lambda_{\rm abs}.
\end{equation}
Compose the unique data-determined incidence congruences around the cycle in
the starting channel frame and denote the resulting Euclidean motion by
\(H_c\).

\begin{lemma}[Periodic cycle-holonomy identity]
\label{lem:v23-cycle-holonomy}
For realizable intrinsic data, \(H_c\) has identity rotational part.  Hence
\(H_c=\tau_{v_c}\) for a data-determined vector \(v_c\).  For every admissible
periodic gluing and every placement \(A_{e_0}(x)=R_{e_0}x+t_{e_0}\) of the
starting channel frame,
\begin{equation}\label{eq:v23-cycle-holonomy}
                         R_{e_0}v_c=\iota(\eta_c).
\end{equation}
Equivalently, after using the global gauge to set the starting frame rotation
to the identity, \(\iota(\eta_c)=v_c\).
\end{lemma}

\begin{proof}
For two consecutive incidences of the same quotient obstacle, the gluing
equation identifies their placed curve copies after the corresponding deck
translations.  The unique signature-rigid congruence between the two recovered
edge-frame copies is therefore conjugate, by their placements, to translation
by the difference of their deck offsets.  Compose these identities around the
cycle.  All intermediate edge placements telescope.  The signed deck
differences sum to \(\eta_c\), giving
\[
 A_{e_0}H_cA_{e_0}^{-1}
 =\tau_{\iota(\eta_c)}.
\]
Conjugation by \(A_{e_0}\) preserves the rotational part, so \(H_c\) is a
translation.  Comparing translation vectors gives
\eqref{eq:v23-cycle-holonomy}.
\end{proof}

\begin{definition}[Lattice-anchoring cycle]
\label{def:v23-lattice-anchoring-cycle}
A signature-rigid cycle is \emph{lattice-anchoring} if \(\eta_c\ne0\).
Because the marked basis fixes the orientation and Gram form of the abstract
lattice, the image of one nonzero vector determines the
orientation-preserving realization \(\iota\) relative to the starting channel
frame through Lemma~\ref{lem:v23-cycle-holonomy}.
\end{definition}

\begin{definition}[Rooted signature-rigid spanning tree]
\label{def:v23-signature-rigid-tree}
Fix a quotient obstacle orbit \(r_*\).  A spanning tree rooted at \(r_*\) is
\emph{signature-rigid} if its edges can be ordered so that each edge joins an
already reached obstacle orbit to one new orbit, and its contact on the already
reached orbit has a curvature signature occurring at a unique framed point of
that obstacle.  For a one-vertex quotient graph the empty tree is allowed.
\end{definition}

\begin{corollary}[Fully intrinsic periodic table rigidity]
\label{cor:v23-signature-rigid-global}
Under the hypotheses of Theorem~\ref{thm:v23-intrinsic-table-rigidity}, assume
the measured network contains a lattice-anchoring cycle and, rooted at an
obstacle orbit met by that cycle, a signature-rigid spanning tree.  Then
\(\mathfrak G_{\rm per}(\mathscr D)\) is a singleton.  Hence the intrinsic
signed lifted-channel data determine the complete labelled periodic billiard
table up to one simultaneous global element of \(\mathrm{SE}(2)\).  No
absolute contact point, absolute tangent frame, inter-channel placement, or
lattice orientation relative to a channel frame is supplied in advance.
\end{corollary}

\begin{proof}
Use the global Euclidean gauge to place the starting frame of the anchoring
cycle.  The cycle places at least one recovered obstacle lift and therefore a
root obstacle image.  Since \(\eta_c\ne0\),
Lemma~\ref{lem:v23-cycle-holonomy} fixes the orientation-preserving lattice
realization \(\iota\) relative to that frame.  Traverse the rooted spanning
tree.  At each edge, the unique signature on the already placed obstacle fixes
the edge-frame transform, and the known deck correction
\(\iota(\ell_e)\) fixes the new quotient obstacle image.  Induction places all
obstacle orbits.  Thus there is at most one algebraic gluing modulo the initial
global gauge.  Since the hypotheses concern data realized by a dispersing
table, that unique gluing is admissible.  Hence the admissible gluing space is
a singleton.
\end{proof}

\begin{corollary}[One-root registration is sufficient once the lattice is known]
\label{cor:v23-known-lattice-frame}
Suppose an oriented Euclidean realization of the periodic lattice and the
Euclidean placement of one root channel frame relative to that lattice are
supplied as ambient calibration.  Then a rooted signature-rigid spanning tree
determines every obstacle image and every remaining measured channel frame.
Thus one root registration suffices; a common laboratory registration of all
channels is not required.
\end{corollary}

\begin{remark}[Symmetry, connectivity, holonomy and admissibility]
\label{rem:v23-generic-holonomy}
The sufficient conditions above are not necessary.  Finite obstacle symmetries
produce finitely many possible incidence matches, and several cycles can
eliminate those choices.  The exact object controlling uniqueness is the
admissible periodic gluing space \(\mathfrak G_{\rm per}(\mathscr D)\).  It
separates four genuine issues: obstacle symmetry, channel-network connectivity,
periodic lattice holonomy and global geometric admissibility.
\end{remark}

\begin{corollary}[Common laboratory registration as a special case]
\label{cor:v23-registered-special-case}
If the Euclidean placement of the lattice and of every measured channel frame
is supplied externally, then the periodic gluing is fixed a priori.  A family
of measured channels meeting every obstacle orbit therefore determines the
whole analytic table in that registration, recovering the registered v22
statement as a special case.
\end{corollary}